% Appendix-local commands/environments (safe when this file is \input into main.tex)
\newcommand{\R}{\mathbb{R}}
\newcommand{\E}{\mathbb{E}}
\newcommand{\1}{\mathbf{1}}

\makeatletter
\@ifundefined{c@remark}{%
  \theoremstyle{remark}%
  \newtheorem{remark}{Remark}%
}{}
\makeatother


\chapter{A Canonical High-Dimensional Model for Averaging-Based Spike Removal and Semantic Preservation}
\label{app:canonical_model}

This appendix provides a detailed, self-contained mathematical instantiation of the
assumptions used in Section~1.6.
The goal is to rigorously justify two claims:
\begin{enumerate}
    \item averaging recovers the shared grammar-driven spike; and
    \item removing this spike preserves essentially all semantic variability.
\end{enumerate}
The model is deliberately high-dimensional and distributional, reflecting the fact
that semantic information in language models is spread over a wide feature space and
manifests through variance rather than mean structure.

\section{Matrix-as-Tokens Model}

Let $G \in \R^{n\times n}$ be a matrix whose columns correspond to $n$ tokens, each
represented by an $n$-dimensional vector.
Denote the $j$-th column by $g_j \in \R^n$, so
\[
G = [g_1,\dots,g_n].
\]

\subsection{Grammar and Semantic Feature Decomposition}

We assume that each token representation decomposes additively as
\begin{equation}
\label{eq:col_decomp}
g_j = g_{\mathrm{gram}} + s_j,
\qquad j=1,\dots,n,
\end{equation}
where:
\begin{itemize}
\item $g_{\mathrm{gram}} \in \R^n$ is a deterministic vector capturing grammatical or
      structural features shared across all tokens; and
\item $s_j \in \R^n$ is a random vector capturing token-specific semantic information.
\end{itemize}

Stacking the semantic components gives
\[
S \triangleq [s_1,\dots,s_n] \in \R^{n\times n},
\]
and therefore
\begin{equation}
\label{eq:G_decomp}
G = g_{\mathrm{gram}}\1^\top + S.
\end{equation}

\section{Semantic Feature Distribution: Wide but Sparse}

We now formalize the notion that semantics are high-dimensional but sparse.

\begin{assumption}[Wide sparse semantic feature model]
\label{assump:sparse_semantics}
There exists a semantic feature dictionary of size $n$.
Each token $j$ activates a random subset $\mathcal{S}_j \subset \{1,\dots,n\}$ with
$|\mathcal{S}_j| = k \ll n$.
The semantic vector takes the form
\[
s_j = \sum_{i \in \mathcal{S}_j} \xi_{j,i} e_i,
\]
where:
\begin{itemize}
\item $\{e_i\}$ is the standard basis of $\R^n$;
\item $\xi_{j,i}$ are i.i.d.\ random variables with
      $\E[\xi_{j,i}] = 0$ and $\E[\xi_{j,i}^2] = \tau^2 < \infty$;
\item the supports $\mathcal{S}_j$ are drawn independently across tokens.
\end{itemize}
\end{assumption}

\begin{remark}
This model captures the intuition that each token expresses only a small subset of
possible semantic features (entities, attributes, relations), while the full semantic
space is very large.
\end{remark}

\begin{assumption}[Finite semantic energy]
\label{assump:energy}
Semantic vectors satisfy
\[
\E[s_j] = 0,
\qquad
\E\|s_j\|_2^2 = \sigma^2 = k\tau^2.
\]
\end{assumption}

\section{Step 1: Averaging Recovers the Grammar Direction}

Define the column mean
\begin{equation}
\label{eq:col_mean}
\bar g \triangleq \frac{1}{n}G\1
= g_{\mathrm{gram}} + \frac{1}{n}\sum_{j=1}^n s_j.
\end{equation}

\begin{proposition}[Column averaging recovers grammar]
\label{prop:mean_recovers_grammar}
Under Assumptions~\ref{assump:sparse_semantics}--\ref{assump:energy},
\[
\bar g \xrightarrow[]{a.s.} g_{\mathrm{gram}}
\qquad \text{as } n\to\infty.
\]
\end{proposition}

\begin{proof}
Fix a coordinate $i\in\{1,\dots,n\}$.
The $i$-th coordinate of $\frac{1}{n}\sum_{j=1}^n s_j$ is
\[
\frac{1}{n}\sum_{j=1}^n s_{j,i}.
\]
Each $s_{j,i}$ is nonzero with probability $k/n$, has zero mean, and finite variance.
By independence across $j$ and Kolmogorov’s strong law of large numbers,
\[
\frac{1}{n}\sum_{j=1}^n s_{j,i} \to 0 \quad \text{almost surely}.
\]
Since this holds for every coordinate and $n$ is finite at each stage,
$\frac{1}{n}\sum_j s_j \to 0$ almost surely as a vector.
\end{proof}

\section{Step 2: Spike Removal via Averaging-Based Projection}

Define the normalized average direction
\[
\hat g \triangleq \frac{\bar g}{\|\bar g\|_2},
\]
and the rank-one projector
\begin{equation}
\label{eq:P_def}
P \triangleq \hat g \hat g^\top,
\qquad
\widetilde G \triangleq (I-P)G.
\end{equation}

This operation removes the shared component using only a matrix--vector product and
normalization, without computing any singular vectors.

\section{Step 3: What Is Removed and What Is Preserved}

\subsection{Grammar Spike Removal}

\begin{proposition}
\label{prop:remove_grammar}
$\|(I-P)g_{\mathrm{gram}}\|_2 \to 0$ almost surely as $n\to\infty$.
\end{proposition}

\begin{proof}
By Proposition~\ref{prop:mean_recovers_grammar}, $\hat g \to g_{\mathrm{gram}}/\|g_{\mathrm{gram}}\|$.
Let $u = g_{\mathrm{gram}}/\|g_{\mathrm{gram}}\|$.
Then $P \to uu^\top$, and
\[
(I-P)g_{\mathrm{gram}} \to (I-uu^\top)g_{\mathrm{gram}} = 0.
\]
\end{proof}

\subsection{Semantic Energy Preservation}

We now show that spike removal eliminates only a vanishing fraction of semantic
energy.

\begin{theorem}[Semantic variance preservation under sparse semantics]
\label{thm:semantic_energy_preserved}
Under Assumptions~\ref{assump:sparse_semantics}--\ref{assump:energy},
\[
\E\|(I-P)S\|_F^2
=
\E\|S\|_F^2 - O(k),
\]
and therefore
\[
\frac{\E\|(I-P)S\|_F^2}{\E\|S\|_F^2}
=
1 - O\!\left(\frac{k}{n}\right).
\]
\end{theorem}

\begin{proof}
For any unit vector $u$, the projection identity gives
\[
\|(I-uu^\top)S\|_F^2 = \|S\|_F^2 - \|u^\top S\|_2^2.
\]
We estimate $\E\|u^\top S\|_2^2$.
Write
\[u^\top S = [u^\top s_1,\dots,u^\top s_n],
\qquad
\|u^\top S\|_2^2 = \sum_{j=1}^n (u^\top s_j)^2.
\]
Taking expectation,
\[
\E\|u^\top S\|_2^2
=
\sum_{j=1}^n \E(u^\top s_j)^2
=
n\, \E(u^\top s_1)^2.
\]

Since $s_1$ has exactly $k$ nonzero coordinates, each contributing $\tau^2$ variance,
and $u$ is dense with $\|u\|_2=1$, we have
\[
\E(u^\top s_1)^2
=
\E\sum_{i\in\mathcal{S}_1} u_i^2 \xi_{1,i}^2
=
\tau^2 \E\sum_{i\in\mathcal{S}_1} u_i^2
=
\tau^2 \cdot \frac{k}{n}.
\]
Therefore
\[
\E\|u^\top S\|_2^2 = n \cdot \tau^2 \cdot \frac{k}{n} = k\tau^2 = O(k).
\]

On the other hand,
\[
\E\|S\|_F^2 = \sum_{j=1}^n \E\|s_j\|_2^2 = n\sigma^2 = nk\tau^2.
\]
Combining,
\[
\E\|(I-P)S\|_F^2
=
nk\tau^2 - k\tau^2
=
k\tau^2(n-1),
\]
and dividing by $\E\|S\|_F^2$ yields the claimed ratio.
\end{proof}

\begin{corollary}
Each token loses only an $O(k/n)$ fraction of its semantic energy under spike removal.
\end{corollary}

\section{Interpretation and Connection to Learning}

This appendix rigorously establishes that:
\begin{itemize}
\item averaging recovers the shared grammar-driven direction (first-order structure);
\item spike removal subtracts only a single shared direction;
\item high-dimensional, sparse semantic variance is preserved up to vanishing
      relative error.
\end{itemize}

Since stochastic gradient descent depends on second-order variability of token-level
updates rather than their mean, preserving semantic variance is sufficient to
maintain semantic learnability.
Thus, the averaging-based spike removal isolates dominant grammatical agreement
structure while leaving semantic learning intact.